\section{Background and Related Work}\label{sec:related}

\subsection{Decentralized Networks}

Since Chaum's seminal work on untraceable communication in 1981~\cite{chaum-mix}, mixnets have been extensively studied as a mechanism for anonymous communication~\cite{piotrowska2017loopix, van2015vuvuzela, kwon2020xrd, lazar2018karaoke, cottrell1995mixmaster, alexopoulos2017MCMIX, chaum2016cmix, chaum-mix, danezis2003mixminion}.
A mixnet is a network of servers, known as \emph{mixnodes}, that forward, delay, and shuffle messages to prevent adversaries from correlating input and output traffic. 
Each packet traverses a sequence of mixnodes before reaching its final destination, as shown in Figure~\ref{fig:mixnet}.

By combining traffic mixing with layered cryptographic encryption, mixnets can provide strong communication anonymity even against powerful traffic-analysis adversaries.
Messages are divided into fixed-size packets and routing information is encrypted in multiple layers, ensuring that each mixnode learns only its immediate predecessor and successor.

A cryptographic packet format is therefore an essential component of a mixnet, as it protects routing information and prevents packets from being tampered or linked across hops. 
Modern mixnet designs build on these principles rely on the Sphinx packet format~\cite{sphinx}.

\begin{figure}%[t]
	\hspace{-9mm}
	\resizebox{1.085\linewidth}{!}{
\newcommand{\figonetextcol}[0]{\color{black}}

\tikzset{  
    innerLink/.style={-, densely dotted, shorten <= 1pt, shorten >= 2pt},  
    outerLink/.style={->, shorten <= 2pt, shorten >= 4pt, >=Latex},
    packet/.style={circle, inner sep=0.2pt, fill=#1},
    A/.style={packet=\colorA!40, font={#1}},
    B/.style={packet=\colorB!50, font={#1}},
    A1/.style={A=1}, A2/.style={A=2}, A3/.style={A=3}, A4/.style={A=4}, A5/.style={A=5}, A6/.style={A=6}, A7/.style={A=7}, A8/.style={A=8}, A9/.style={A=9},
    B1/.style={B=1}, B2/.style={B=2}, B3/.style={B=3}, B4/.style={B=4}, B5/.style={B=5}, B6/.style={B=6}, B7/.style={B=7}, B8/.style={B=8}, B9/.style={B=9},
}

\newcommand{\countitems}[1]{%
  \def\size{0}%
  \foreach \x in {#1} {%
    \pgfmathparse{\size+1}%
    \global\let\size=\pgfmathresult%
  }%
}

\newcommand{\outerLink}[3]{
  % size (nbr of packets)
  \countitems{#3} 
  % If odd nbr of packet, slight shift to not have a packet on an intersection
  \pgfmathparse{mod(\size,2)==1 ? \size-0.5 : \size}%
  \xdef\size{\pgfmathresult}%

  % define line
  \coordinate (start) at (#1.east);
  \coordinate (stop) at (#2.west);
  \draw[outerLink] (start) -- (stop);
  % put packet on the line
  \foreach \packet [count=\i from 1] in {#3} {
    \pgfmathsetmacro{\t}{1 - \i/(\size+1)}%
    \path let
      \p1 = (start),
      \p2 = (stop),
      \p3 = ($(\p1)!\t!(\p2)$)
    in
      node[\packet] at (\p3) {};
  }
}


\newcommand{\mixnode}[2]{%
  \countitems{#1}%

  \begin{tikzpicture}
    \def\d{0.8}
    \draw[thick] (0,0) circle (\d);

    % Connect according to mapping
    \pgfmathsetmacro{\xoffset}{(1.14*\d - \size/10)*\d}
    \foreach \r [count=\l from 1] in {#1} {
      \pgfmathsetmacro{\yposL}{\d - \l/(\size+1)*2*\d}
      \coordinate (left\l) at (-\xoffset, \yposL);
      \pgfmathsetmacro{\yposR}{\d - \r/(\size+1)*2*\d}
      \coordinate (right\l) at (\xoffset, \yposR);
    }
    % Draw dots
    \foreach \packet [count=\i from 1] in {#2} {
      \draw[innerLink] (left\i) -- (right\i);
      \node[\packet] at (left\i) {};
      \node[\packet] at (right\i) {};
    }
  \end{tikzpicture}
}

\newcommand{\client}[1]{
  \begin{tikzpicture}[scale=0.4]
    % Monitor screen
    \draw[fill=#1!20, draw=black, rounded corners=4pt] (0,0) rectangle (4,2.5);
    % Inner bezel
    \draw[fill=white, draw=none, rounded corners=3pt] (0.2,0.2) rectangle (3.8,2.3);
    % Stand
    \draw[fill=#1!40, draw=black] (1.75,-0.7) rectangle (2.25,0);
    % Base
    \draw[fill=#1!40, draw=black, rounded corners=2pt] (0.8,-1) rectangle (3.2,-0.7);
  \end{tikzpicture}
}

\newcommand{\server}[1]{
  \begin{tikzpicture}[scale=0.7]
    \foreach \y in {0, -0.6, -1.2} {
        \draw[fill=#1!20] (0,\y) rectangle (2,\y-0.5); % rectangle
        \draw (0.2,\y-0.25) circle (0.05); % power/status LED
    }
  \end{tikzpicture}
}

\newcommand{\colorA}{blue}
\newcommand{\colorB}{orange}

\newcommand{\clientA}{\client{\colorA}}
\newcommand{\clientB}{\client{\colorB}}
\newcommand{\serverA}{\server{\colorA}}
\newcommand{\serverB}{\server{\colorB}}


\begin{tikzpicture}
  \def\dx{4}
  \def\dy{-1}

  % GENERAL ARCHITECTURE (define nodes)
  \node (clientA) at (0,-\dy) {\clientA};
  \node (clientB) at (0,\dy) {\clientB};

  \node (serverA) at (4*\dx,-\dy) {\serverA};
  \node (serverB) at (4*\dx,\dy) {\serverB};

  % PACKETS - links and shuffles (python generated)
\node (mix11) at (1*\dx,-\dy) {\mixnode{3,1,2}{A2,A1,B3}};
\node (mix12) at (1*\dx, \dy) {\mixnode{3,2,1}{B2,B1,A3}};
\node (mix21) at (2*\dx,-\dy) {\mixnode{3,1,2}{A3,A2,B3}};
\node (mix22) at (2*\dx, \dy) {\mixnode{2,3,1}{B2,B1,A1}};
\node (mix31) at (3*\dx,-\dy) {\mixnode{3,1,4,2}{B2,A2,A1,B3}};
\node (mix32) at (3*\dx, \dy) {\mixnode{2,1}{A3,B1}};

\outerLink{clientA}{mix11}{A1,A2};
\outerLink{clientA}{mix12}{A3};
\outerLink{clientB}{mix11}{B3};
\outerLink{clientB}{mix12}{B1,B2};

\outerLink{mix11}{mix21}{A2,B3};
\outerLink{mix11}{mix22}{A1};
\outerLink{mix12}{mix21}{A3};
\outerLink{mix12}{mix22}{B2,B1};

\outerLink{mix21}{mix31}{A2,B3};
\outerLink{mix21}{mix32}{A3};
\outerLink{mix22}{mix31}{B2,A1};
\outerLink{mix22}{mix32}{B1};

\outerLink{mix31}{serverA}{A2,A1};
\outerLink{mix31}{serverB}{B2,B3};
\outerLink{mix32}{serverA}{A3};
\outerLink{mix32}{serverB}{B1};

  % ANNOTATIONS
  \node at (0*\dx, -2.2*\dy) {\figonetextcol Clients};
  \draw[loosely dotted, fill=gray!30] (0.23*\dx, -2.2*\dy) -- (0.23*\dx, +2*\dy) {};
  \node at (0.5*\dx, -2.2*\dy) {\figonetextcol $t_0$};
  \draw[loosely dotted, fill=gray!30] (0.77*\dx, -2.2*\dy) -- (0.77*\dx, +2*\dy) {};
  
  \node at (1*\dx, -2.2*\dy) {\figonetextcol $t_1$};
  \draw[loosely dotted, fill=gray!30] (1.23*\dx, -2.2*\dy) -- (1.23*\dx, +2*\dy) {};
  \node at (1.5*\dx, -2.2*\dy) {\figonetextcol $t_2$};
  \draw[loosely dotted, fill=gray!30] (1.77*\dx, -2.2*\dy) -- (1.77*\dx, +2*\dy) {};
  
  \node at (2*\dx, -2.2*\dy) {\figonetextcol $t_3$};
  \draw[loosely dotted, fill=gray!30] (2.23*\dx, -2.2*\dy) -- (2.23*\dx, +2*\dy) {};
  \node at (2.5*\dx, -2.2*\dy) {\figonetextcol $t_4$};
  \draw[loosely dotted, fill=gray!30] (2.77*\dx, -2.2*\dy) -- (2.77*\dx, +2*\dy) {};
  
  \node at (3*\dx, -2.2*\dy) {\figonetextcol $t_5$};
  \draw[loosely dotted, fill=gray!30] (3.23*\dx, -2.2*\dy) -- (3.23*\dx, +2*\dy) {};
  \node at (3.5*\dx, -2.2*\dy) {\figonetextcol $t_6$};
  \draw[loosely dotted, fill=gray!30] (3.77*\dx, -2.2*\dy) -- (3.77*\dx, +2*\dy) {};
  \node at (4*\dx, -2.2*\dy) {\figonetextcol Services};

\end{tikzpicture}
}
	\caption{Step-by-step tracing of two clients communicating with services through a 3-layer mixnet with two mixnodes (large circles) at each layer.}
	\label{fig:mixnet}
\end{figure}

\subsection{Sphinx Packet Format}

The \emph{Sphinx} packet format~\cite{sphinx} provides three essential security properties: per-hop confidentiality, per-hop integrity protection, and strong unlinkability.

First, Sphinx ensures \emph{per-hop confidentiality} by encrypting routing information in layers, similar to classical onion routing.  
Each mixnode learns only its immediate predecessor and successor, but gains no knowledge about the overall path, the sender, or the final destination.  

Second, Sphinx provides \emph{per-hop integrity protection} through a cryptographic tag embedded in each layer of the header. 
Any modification results in an integrity-tag verification failure at the mixnode, causing the packet to be discarded.
This prevent adversaries from tampering with routing information.  

Third, Sphinx offers \emph{strong unlinkability} by ensuring that packets entering and leaving a mixnode are computationally indistinguishable.  
This is achieved by using fixed-size headers and payloads, and by re-randomizing cryptographic elements at every hop.
Each hop transforms the packet into a fresh, uniformly distributed header, preventing an adversary from linking an incoming packet to its outgoing counterpart.

These guarantees have made Sphinx the de facto standard for privacy-preserving routing in modern mixnet-based systems like Nym~\cite{nym} and Katzenpost~\cite{infeld2025echomix}.  
Despite all properties provided by the Sphinx, it does not solve one of the anonymous communication challenge which is privacy-preserving acces control.
To solve this, Sphinx can be combined with anonymous credentials.

\subsection{Anonymous Credentials}

Credential is a set of personal attributes bound to its owner by cryptographic means that are signed, and thereby certified, by the issuer (i.e. certifying party) \cite{Fisher_PET}.
Credentials provide a mechanism for access-control by proving authorization or possession of certified attributes.
Anonymous credentials are rerandomizable credentials which are able to demonstrate eligibility without revealing the identity of the credential holder.

Anonymous credentials also often offers selective-disclosure property which allows proving only a subset of its attributes to a verifier and keep the remaining attributes hidden.
In addition, instead of revealing the exact value of an attribute, anonymous credential can also enable the user to reveal only attribute property without revealing the attribute itself (e.g. age > 18, instead of age = 26).
Therefore, these properties makes anonymous credentials particularly suitable for privacy-preserving access-control systems.

Coconut~\cite{coconut} is a threshold issuance anonymous credential with selective disclosure designed for decentralized systems. 
It combines selective disclosure credentials with threshold issuance, allowing a set of authorities to jointly issue credentials without requiring a single authority to hold the complete signing key. 
A client can obtain partial credentials from multiple authorities and combine them into a valid credential, while the use of blinding and re-randomization prevents the authorities from linking credential issuance to subsequent presentations. 
This distributed-trust model has made Coconut a useful building block for privacy-preserving access-control in mixnet deployments such as Nym.

% HERE IS THE PROBLEM: proof of use
Although Coconut credentials provide privacy-preserving authorization, they primarily offer a \emph{proof of possession}: a client can prove that it is authorized for a given destination without revealing its identity. 
However, they do not ensure that the authorization is actually being used for the destination encoded inside the Sphinx packet. 
Since the destination remains encrypted as the packet traverses the mixnet, a client could potentially present a credential for an authorized destination while constructing the packet to reach a different, unauthorized destination.
This deception can only be detected at the exit gateway, after the packet has already consumed network resources. 
This motivates extending the Sphinx packet format with a credential that is cryptographically bound to the encrypted destination. 
Our construction allows mixnodes to verify that a packet carries valid authorization for its destination 
without learning either the destination or the client's identity, while preserving the privacy properties of anonymous credentials.
